\documentclass[12pt]{article}

\usepackage[margin=1in]{geometry}
\usepackage{lmodern}
\usepackage[T1]{fontenc}
\usepackage{amsmath, amssymb, amsthm}
\newtheorem{theorem}{Theorem}
\newtheorem{definition}{Definition}
\newtheorem{proposition}{Proposition}
\newtheorem{remark}{Remark}
\usepackage{graphicx}
\usepackage{hyperref}
\usepackage{natbib}
\usepackage{booktabs}
\usepackage{xcolor}
\usepackage{multirow}
\usepackage{array}
\usepackage{microtype}
\usepackage{setspace}
\usepackage{titlesec}
\usepackage{footmisc}
\usepackage{mdframed}

%% Section formatting
\titleformat{\section}{\normalfont\large\bfseries}{\thesection.}{0.5em}{}
\titleformat{\subsection}{\normalfont\normalsize\bfseries}{\thesubsection.}{0.5em}{}
\titleformat{\subsubsection}{\normalfont\normalsize\itshape}{\thesubsubsection.}{0.5em}{}

\setstretch{1.4}

%% Shorthand macros
\newcommand{\ar}{\textsc{ApportionRegions}}
\newcommand{\recom}{\textsc{ReCom}}
\newcommand{\pp}{\ensuremath{\mathrm{PP}}}

%% Legal case macros
\newcommand{\case}[1]{\textit{#1}}
\newcommand{\rucho}{\case{Rucho v.\ Common Cause}}
\newcommand{\callais}{\case{Louisiana v.\ Callais}}

%% Note markers
\newcommand{\est}{$^{\dagger}$}

%% Highlighted box
\newenvironment{keybox}{%
  \begin{mdframed}[linewidth=0.8pt, innerleftmargin=12pt, innerrightmargin=12pt,
                   innertopmargin=8pt, innerbottommargin=8pt,
                   backgroundcolor=gray!8]
  \small\setstretch{1.3}%
}{%
  \end{mdframed}%
}

%% -----------------------------------------------------------------------

\title{%
  \textbf{Are Algorithmic Maps Partisan?}\\
  \textbf{Situating a BISECT Benchmark Within ReCom Ensembles}\\[0.4em]
  \large Two Implementations and a Complete 2020 Election Model\\[0.8em]
  \normalsize\textit{G.2 --- Partisan Outcome Distributions}%
}

\author{%
  Gio Della-Libera\\
  \textit{Apportionment Research Group}\\
  \small\texttt{giodl@microsoft.com}
}

\date{May 2026}

\begin{document}

\maketitle

\begin{abstract}
\noindent
We report real tract-level partisan-seat distributions for Rhode Island, Iowa,
and North Carolina using the 2020 presidential estimates and two ReCom
implementations. Rhode Island has no seat variation: the two-seat benchmark and
all sampled plans elect two Democratic-majority districts under this model.
Iowa's one-Democratic-seat benchmark falls at the 34th--37th percentile.
North Carolina's five-seat benchmark falls at the 17th--21st percentile, but
ESS is below 100 for both implementations and the result is not a converged
tail finding. The locally available 2016 estimates fail current-tract coverage
and are excluded from the primary results. These findings do not support the
former six-state ``near median'' generalization.
\end{abstract}

\tableofcontents
\newpage

\section{Election Model And Coverage}

The primary metric sums tract-level Democratic and Republican presidential
estimates within each district and counts Democratic-plurality districts. The
2020 file matched every Rhode Island, Iowa, and North Carolina benchmark tract.
This is a descriptive seats model, not a legal fairness test
~\citep{stephanopoulos2015partisan,rodden2019why}.
The 2016 file missed 12 of 250 Rhode Island tracts, 129 of 896 Iowa tracts, and
942 of 2,672 North Carolina tracts because it does not share the 2020 tract
universe. No crosswalk was available, so 2016 is excluded from the primary
percentiles rather than treating unmatched units as zero votes.

\section{2020 Presidential Seat Results}

\begin{table}[h]
\centering
\caption{Benchmark Democratic-seat position under the 2020 presidential tract model.}
\begin{tabular}{lrrrrrr}
\toprule
State & Tool & Benchmark & Mean & Percentile & $\widehat R$ & ESS \\
\midrule
RI & Rust & 2 & 2.000 & 50.0 & 1.000 & 6000 \\
RI & GerryChain & 2 & 2.000 & 50.0 & 1.000 & 6000 \\
IA & Rust & 1 & 1.266 & 36.7 & 1.002 & 486 \\
IA & GerryChain & 1 & 1.323 & 33.9 & 1.006 & 1072 \\
NC & Rust & 5 & 5.944 & 17.1 & 1.028 & 90 \\
NC & GerryChain & 5 & 5.811 & 21.2 & 1.036 & 98 \\
\bottomrule
\end{tabular}
\end{table}

Rhode Island is degenerate for this metric: every sampled plan has two
Democratic seats, so the mid-rank percentile is 50\% but carries no evidence of
partisan neutrality. Iowa is below the ensemble mean but not in an extreme
tail. North Carolina is lower than most sampled plans, but both ESS values fail
the preregistered 100-sample threshold; the 17--21\% estimates remain
non-converged descriptive results.

\section{Interpretation}

The state results are heterogeneous and election-specific. They do not show
that the benchmark systematically favors or disfavors a party. They also do
not support the earlier claim that the algorithm is near the partisan median
in five of six states. Cross-tool seat means are similar here, while the
underlying cut distributions differ materially.

The proportionality corridor is not reported because the experiment did not
preregister a statewide proportional-entitlement rule. A seat percentile under
one presidential election is a diagnostic, not a fairness standard.

\section{Claim Boundary}

The package is tract-level, covers three states and one complete election, and
does not establish intent, causation, partisan fairness, or legal validity.
Additional elections require a valid crosswalk to the same unit universe.


\bibliographystyle{plainnat}
\bibliography{references}

\end{document}
